% %%%%%%%%%%%%%%%%%%%%%%%%%%%%%%%%%%%%%%%%%%%%%%%%%%%%%%%%%%%%%%%%%%%%%
% %% State-of-the-Art %%
% Gamma
% %%%%%%%%%%%%%%%%%%%%%%%%%%%%%%%%%%%%%%%%%%%%%%%%%%%%%%%%%%%%%%%%%%%%%

\subsection{Gamma}
\label{sec:soa:gamma}

Gamma~\cite{Zhang2021_Gamma} is a dedicated accelerator for SpMSpM kernels, which uses \algo{Gust} algorithm. The two input matrices are both stored in \emph{(UOP, CP) row-major} format and the output is generated in the same format, providing consistency.

% Figure in Accelerators.tex

The Gamma architecture, presented in Figure~\ref{fig:soa:gamma}, is composed of a fetcher for the $A$-rows that contains a \emph{Distributor} to distribute them to 32 \emph{PEs}. Each \emph{PE} processes an entire $A$-row and computes the corresponding $C$-row. The \emph{PE} fetches the required $B$-rows and merges and sorts them in a 64-wide \emph{merge} block, before multiplying with the $A$-elements. As a result, the output $C$-row  is generated in order. The \emph{reduction} step is, thus, highly simplified, because any duplicate indices are adjacent. Normally, the $C$-row can be written back to main memory. However, if the $A$-row contains more than 64 non-zero elements, then it must be processed with multiple passes. After each pass, the partial $C$-row is stored temporarily. Finally, the temporary $C$-rows are read-back and merged by the \emph{PEs}.
Of course, the $B$-rows must be accessed repeatedly by the different \emph{PEs}, and Gamma proposes \emph{FiberCache} which is a highly banked cache which fetches the $B$-rows in advance, based on the column index table of $A$. It keeps track of which rows of $B$ are most needed by the \emph{PEs}, and when an eviction is necessary, the victim, is the row which is least needed.

Gamma encounters some difficulties when the $A$-matrix has rows that are too dense, requiring multiple iterations through the \emph{PEs}, which creates a high-load on the cache. The authors propose a software preprocessing technique which splits dense $A$-rows and rearranges them to make best use of the \emph{FiberCache}.